\section{KNN Classification}

In this section, we apply the K-Nearest Neighbors (KNN) algorithm to classify patients based on their
clinical measurements. KNN is an instance-based, non-parametric learning method and is particularly
appropriate for this dataset because no strong linear separability was observed in the EDA, and several
features (especially \texttt{thalach} and \texttt{oldpeak}) exhibit nonlinear relationships with the
presence of heart disease. Since KNN relies directly on the distance between feature vectors, proper
preprocessing and feature scaling are crucial components of the classification pipeline.

\subsection{Overview of the KNN Algorithm}

Given a test instance $x$, the KNN classifier operates by:

\begin{enumerate}
	\item Computing the distance between $x$ and all training samples,
	\item Selecting the $k$ closest training samples,
	\item Determining the majority class among these $k$ neighbors,
	\item Assigning that class label to $x$.
\end{enumerate}

Because KNN does not assume any particular probabilistic or functional form for the data, it is well suited
for medical datasets where complex and nonlinear interactions exist among variables---as confirmed by the
pair plots and correlation patterns observed in our EDA.

\subsection{Mixed Distance Function for Numeric and Categorical Features}

The dataset contains a mixture of numeric variables (such as \texttt{age}, \texttt{trestbps},
\texttt{chol}, \texttt{thalach}, and \texttt{oldpeak}) and categorical or binary attributes
(e.g., \texttt{sex}, \texttt{cp}, \texttt{exang}, \texttt{restecg}). Using a single purely
numeric distance metric is therefore inappropriate.

To address this, a mixed distance measure combining Euclidean distance for numeric features and Hamming
distance for categorical features was implemented:

\[
d(x,y) = w_{\text{num}} \cdot d_{\text{num}}(x,y) +
w_{\text{cat}} \cdot d_{\text{cat}}(x,y),
\]

where both weights were set to $0.5$ in this study. This weighting ensures that numeric attributes—
particularly influential ones such as \texttt{thalach} and \texttt{oldpeak}—contribute proportionally
to the similarity computation, without overwhelming information provided by categorical features.

\subsection{Model Training and Choice of $k$}

After preprocessing and Min--Max normalization, the dataset was split into an 80/20 training–testing partition.
The KNN classifier was trained on the resulting \texttt{X\textsubscript{train}} matrix and evaluated on
\texttt{X\textsubscript{test}}.

To examine the effect of the neighborhood size, we evaluated KNN for:

\[
k \in \{3,\,5,\,7\}.
\]

Performance was assessed using accuracy, precision, recall, and F1-score—metrics that are especially
important in clinical contexts, where both false positives and false negatives carry significant
consequences.

\subsection{Classification Results}

\begin{table}[H]
	\centering
	\caption{Performance of the KNN classifier for different values of $k$.}
	\label{table:knn-results}
	\begin{tabular}{cccccc}
		\toprule
		$k$ & Accuracy & Precision & Recall & F1-score \\
		\midrule
		3 & 0.9317 & 0.9159 & 0.9515 & 0.9333 \\
		5 & 0.8195 & 0.7946 & 0.8641 & 0.8279 \\
		7 & 0.8341 & 0.8053 & 0.8835 & 0.8426 \\
		\bottomrule
	\end{tabular}
\end{table}

The highest accuracy was achieved at $k = 3$. This is consistent with the moderate level of class overlap
observed in the scatter plots and pair plots: smaller values of $k$ allow the classifier to better capture
local patterns driven by the two most informative variables (\texttt{thalach} and \texttt{oldpeak}).
Larger values of $k$ tend to smooth out these local distinctions, causing a drop in performance.

Thus, the optimal neighborhood size for this dataset was:

\[
k_{\text{best}} = 3.
\]

\subsection{Confusion Matrix}

The confusion matrix for the best-performing model is shown in Figure~\ref{fig:confusion-best}.

\begin{figure}[H]
	\centering
	\includegraphics[width=0.6\textwidth]{confusion_matrix.png}
	\caption{Confusion matrix of the KNN classifier for $k=3$.}
	\label{fig:confusion-best}
\end{figure}

The classifier achieves a good balance between sensitivity and specificity. The relatively high recall
indicates that the majority of heart disease cases were correctly identified. A small number of false
positives and false negatives remain, which is expected given the overlapping distributions observed in
the EDA—particularly for features such as \texttt{age} and \texttt{chol}, which show weak class
separation. Nevertheless, the model overall demonstrates strong predictive performance given the mixed
nature of the data and the limited number of highly discriminative features.